% =========================================================
\chapter{Conclusiones}
\label{chapter:conclusiones}
% =========================================================

Este trabajo ha investigado si el balanceo demográfico de los conjuntos
de entrenamiento ---aplicado sobre las dimensiones de edad y sexo---
mejora el rendimiento predictivo de un modelo LSTM de predicción de
glucosa entrenado sobre datos CGM. La pregunta de investigación no admite
una respuesta binaria: los resultados demuestran que el efecto del
balanceo es fuertemente dependiente del repositorio, del rango glucémico
que se priorice y de la técnica específica empleada. Esta sección sintetiza
las conclusiones principales, evalúa el grado de cumplimiento de cada
objetivo específico y propone las líneas de trabajo futuro más prometedoras.

% ---------------------------------------------------------
\section{Conclusiones principales}
\label{sec:conclusiones_principales}
% ---------------------------------------------------------

\paragraph{El balanceo demográfico no mejora el rendimiento global de
forma generalizable.}
En dos de los tres repositorios evaluados ---T1DiabetesGranada y
REPLACE-BG--- el test de Friedman no rechaza la hipótesis nula de
equivalencia entre las 21 condiciones experimentales en ningún rango
glucémico ni métrica de error (RMSE ENTIRE: $p = 0{,}120$ y $p = 0{,}140$
respectivamente). Las diferencias observadas son estadísticamente
compatibles con variación aleatoria, y ninguna técnica de balanceo
supera consistentemente al \textit{baseline} original. Solo en DiaTrend
el balanceo produce efectos estadísticamente significativos ($p < 0{,}0001$
en ENTIRE RMSE), pero incluso allí el efecto es ambivalente: las técnicas
que mejoran el rendimiento global deterioran el rendimiento en hipoglucemia
severa (TBR-2), y viceversa.

\paragraph{Existe un trade-off estructural entre rendimiento global y
sensibilidad a hipoglucemia severa.}
Este es el hallazgo clínicamente más relevante del trabajo. En DiaTrend,
el único repositorio con suficiente heterogeneidad demográfica y tamaño
para generar efectos detectables, ninguna técnica evaluada consigue
mejorar simultáneamente el RMSE en el rango completo (ENTIRE) y en
TBR-2. Las técnicas de aumento (ROS, SMOTE, Jittering) mejoran el RMSE
global al enriquecer la representación de perfiles glucémicos de grupos
minoritarios con mejor control medio, pero a costa de reducir la
exposición relativa a ventanas de hipoglucemia severa durante el
entrenamiento. Las técnicas de submuestreo producen el efecto opuesto.
Este trade-off es coherente con los cambios en la distribución glucémica
del \textit{train} documentados en la sección~\ref{subsec:impacto_glucemico}
y tiene implicaciones directas para el diseño de sistemas de predicción
clínica: la elección de la técnica debe hacerse en función del rango
glucémico que se desea priorizar, con plena conciencia del coste en los
rangos restantes.

\paragraph{El efecto del balanceo es dataset-dependiente y explicable
por propiedades estructurales.}
Los tres repositorios representan tres escenarios cualitativamente
distintos. T1DiabetesGranada, con gran volumen y desequilibrio moderado,
es insensible al balanceo: el modelo ya dispone de suficiente información
con la partición original. REPLACE-BG, con seguimientos estandarizados
que minimizan la heterogeneidad inter-paciente, elimina estructuralmente
el problema que el balanceo intenta corregir. DiaTrend, con desequilibrio
extremo (38:1 por edad) y seguimientos muy heterogéneos, es el único
escenario donde el balanceo produce un efecto real, pero también el
más frágil metodológicamente dado su pequeño tamaño muestral. Esta
dependencia estructural implica que no es posible hacer una recomendación
universal de técnica: la elección óptima debe evaluarse caso a caso en
función de las propiedades del repositorio.

\paragraph{El balanceo por edad supera al balanceo por sexo en magnitud
de efecto.}
La dimensión edad produce variaciones de rendimiento sistemáticamente
más amplias que la dimensión sexo en todos los repositorios, coherente
con su mayor ratio de desequilibrio (hasta 38:1 frente a 2{,}2:1 en
sexo) y su mayor correlación con el perfil glucémico. En DiaTrend las
técnicas de edad producen una dispersión de $\approx$3\,mg/dL en RMSE
(desde $-3\%$ hasta $-7\%$ respecto al \textit{baseline}), mientras que
las de sexo se concentran en $\pm$0{,}5\,mg/dL. En T1DiabetesGranada
y REPLACE-BG ambas dimensiones son estadísticamente equivalentes, pero
las técnicas de edad presentan mayor variabilidad inter-fold. Si los
recursos computacionales son limitados, el balanceo por edad debe
priorizarse.

\paragraph{Tomek Links es la técnica con mejor relación
beneficio/coste en el repositorio más complejo.}
En DiaTrend, Tomek Links (dimensión edad) obtiene el menor RMSE global
de todas las condiciones ($44{,}83$\,mg/dL, posición 1 en el CD diagram
de ENTIRE) sin deteriorar el RMSE en TBR-2, siendo la única técnica que
evita el trade-off estructural en ese repositorio. Su mecanismo ---eliminar
exclusivamente puntos fronterizos ambiguos--- produce una reducción
moderada del volumen de entrenamiento ($-1{,}8\%$ en el grupo $<$31)
que mejora la calidad del aprendizaje sin los riesgos de redundancia del
\textit{oversampling} ni la pérdida de información del \textit{undersampling}
agresivo. En T1DiabetesGranada y REPLACE-BG no produce ni mejoras ni
deterioros apreciables, lo que lo convierte en la opción de menor riesgo
entre todos los escenarios evaluados.

% ---------------------------------------------------------
\section{Grado de cumplimiento de los objetivos}
\label{sec:cumplimiento_objetivos}
% ---------------------------------------------------------

\textbf{Objetivo 1 — Diseño e implementación del \textit{pipeline} de
balanceo.} Cumplido. Se diseñó e implementó un módulo de balanceo que
opera exclusivamente sobre el subconjunto \texttt{'train'} de cada
\textit{fold}, garantizando la invarianza bit a bit de los conjuntos de
validación y test entre todas las condiciones experimentales. El módulo
genera archivos \textit{parquet} intermedios verificables que permiten
auditar el estado exacto del conjunto de entrenamiento sin necesidad de
re-ejecutar el balanceo, y fue integrado en el \textit{pipeline} previo
del proyecto RELIEF-T1D sin modificar su lógica de entrenamiento.

\textbf{Objetivo 2 — Implementación y validación de la taxonomía de
técnicas.} Cumplido. Se implementaron y validaron diez técnicas de
balanceo organizadas en tres familias (aleatorias, guiadas e híbridas),
aplicadas en dos variantes demográficas (edad y sexo) más una condición
de control, resultando en 40 condiciones experimentales en
T1DiabetesGranada y REPLACE-BG y 39 en DiaTrend (por ausencia de la
condición \textit{balanced\_age\_oversampling} en ese repositorio). La
validación de cada artefacto generado comprobó propiedades de
exclusividad mutua, cobertura completa de pacientes, equilibrio de carga
entre \textit{folds} y cumplimiento de las restricciones de tamaño por
familia.

\textbf{Objetivo 3 — Evaluación del efecto sobre el rendimiento del
modelo LSTM.} Cumplido. Se evaluó el rendimiento del modelo LSTM mediante
RMSE y MAE desagregados por rango glucémico (TBR-2, TBR-1, TIR, TAR-1,
TAR-2 y ENTIRE) en los tres repositorios y las 40/39 condiciones
experimentales. El análisis documentó tanto el nivel central de rendimiento
como su variabilidad inter-\textit{fold} y la disociación entre métricas
globales y clínicas.

\textbf{Objetivo 4 — Comparación del impacto por edad frente a sexo.}
Cumplido. El análisis del gráfico de pendientes
(figura~\ref{fig:slope_rmse_dim}), los heatmaps de variación demográfica
y los CD diagrams demostraron que el balanceo por edad produce efectos
más amplios y más variables que el balanceo por sexo en todos los
repositorios, con la brecha más pronunciada en DiaTrend. Se identificaron
las razones estructurales de esta diferencia (mayor ratio de desequilibrio,
mayor correlación entre grupo etario y perfil glucémico, mayor
amplificación al pasar de nivel de paciente a nivel de ventana).

\textbf{Objetivo 5 — Caracterización de la robustez entre datasets.}
Cumplido. Los resultados son marcadamente heterogéneos entre repositorios:
el test de Friedman rechaza la hipótesis nula solo en DiaTrend, y los
patrones de mejora/deterioro difieren cualitativamente entre los tres
conjuntos. Se vincularon estas diferencias con propiedades estructurales
documentadas en el EDA: tamaño muestral, patrón MNAR, ratio de
desequilibrio demográfico y homogeneidad de los seguimientos.

\textbf{Objetivo 6 — Identificación de la técnica con mejor rendimiento
agregado.} Cumplido. El ranking Borda y los CD diagrams
de Nemenyi identifican a Tomek Links (dimensión edad) como la técnica
con mejor rendimiento agregado en DiaTrend ---el único repositorio donde
el balanceo produce efectos estadísticamente significativos--- y a
Age·Jittering como la mejor en T1DiabetesGranada y REPLACE-BG según el
Borda, aunque sin significación estadística. No existe ninguna técnica
que domine en los tres repositorios simultáneamente, lo que impide
formular una recomendación universal. El objetivo se considera cumplido
en la medida en que se identificó la mejor técnica dentro de cada
escenario y se justificó la ausencia de un ganador universal.

\textbf{Objetivo 7 — Descripción de limitaciones y compromisos clínicos.}
Cumplido. La sección~\ref{sec:discusion} documenta el trade-off entre
rendimiento global y sensibilidad a hipoglucemia severa, cuantifica su
magnitud en DiaTrend y describe las cuatro limitaciones metodológicas
principales: potencia estadística limitada por el número de
\textit{folds}, tamaño insuficiente de DiaTrend para generalizaciones
robustas, restricción a un único modelo y horizonte de predicción, y
ausencia de análisis de las propiedades temporales de las series generadas.

% ---------------------------------------------------------
\section{Trabajos futuros}
\label{sec:trabajos_futuros}
% ---------------------------------------------------------

Los resultados de este trabajo abren varias líneas de investigación que
no pudieron abordarse dentro de su alcance.

\textbf{Ampliación del protocolo estadístico.} El principal factor
limitante de este trabajo es la potencia del test de Friedman con solo
5 \textit{folds}. Una extensión natural consiste en repetir los
experimentos con múltiples semillas aleatorias para la asignación de
pacientes a \textit{folds}, generando distribuciones de métricas más
robustas y reduciendo la diferencia crítica de Nemenyi a valores que
permitan detectar efectos de magnitud pequeña. Alternativamente, la
adopción de un esquema de validación cruzada anidada permitiría separar
la selección de técnica de la evaluación de su rendimiento.

\textbf{Extensión a otros modelos y horizontes de predicción.} Los
efectos del balanceo demográfico observados en este trabajo son específicos
del modelo LSTM con horizonte PH$=4$. Sería relevante estudiar si los
mismos patrones se reproducen con modelos basados en transformers o con
enfoques de aprendizaje federado, especialmente dado que estos últimos
son particularmente adecuados para el dominio clínico donde la privacidad
de los datos es crítica. La extensión a horizontes más cortos (PH$=2$,
30 min) o más largos (PH$=8$, 120 min) permitiría determinar si el
trade-off entre RMSE global y TBR-2 varía con el horizonte temporal.

\textbf{Balanceo simultáneo por edad y sexo.} Este trabajo evaluó las
dos dimensiones demográficas de forma independiente. Una extensión
relevante es el balanceo conjunto por la combinación edad$\times$sexo,
que el análisis del heatmap (figura~\ref{fig:age_sex_heatmap}) mostró
que presenta celdas con representación nula o muy baja en algunos
repositorios. Esta extensión requiere técnicas de síntesis condicionada
que respeten la estructura de correlación entre ambas dimensiones,
un problema metodológico abierto en el campo del balanceo demográfico
para series temporales.

\textbf{Análisis del efecto sobre las propiedades temporales de las
series.} Las técnicas de síntesis evaluadas (SMOTE, Jittering) generan
ventanas nuevas modificando los valores glucémicos pero sin garantizar
que las propiedades temporales de las series originales ---autocorrelación,
estacionariedad local, distribución de transiciones glucémicas--- se
preserven en los datos sintéticos. Un análisis de estas propiedades en
los \textit{train sets} balanceados, y su correlación con el rendimiento
del LSTM, permitiría entender mejor el mecanismo por el que el balanceo
afecta al aprendizaje del modelo y diseñar técnicas de síntesis más
adecuadas para datos CGM.

\textbf{Modelos de fundación para series temporales como alternativa al
balanceo.} Una línea emergente que podría cambiar el planteamiento de
este problema es la aparición de modelos de fundación preentrenados
sobre grandes colecciones de series temporales heterogéneas (como
TimesFM, Moirai o Chronos). Estos modelos incorporan durante el
preentrenamiento una diversidad demográfica y clínica mucho mayor que
la disponible en cualquier cohorte individual, lo que plantea la pregunta
de si el balanceo demográfico del conjunto de entrenamiento sigue siendo
necesario cuando el modelo base ya ha aprendido representaciones generales
de la dinámica glucémica a partir de millones de series. Estudiar si el
ajuste fino (\textit{fine-tuning}) de estos modelos sobre cohortes
locales mantiene o elimina los sesgos demográficos documentados en este
trabajo constituye una dirección de investigación relevante y
metodológicamente novedosa en el campo de la predicción CGM.

\textbf{Aprendizaje federado con balanceo demográfico distribuido.}
En entornos clínicos reales, los datos de distintos hospitales o regiones
no pueden centralizarse por razones de privacidad. El aprendizaje
federado permite entrenar modelos colaborativamente sin compartir datos
individuales, pero introduce un nuevo nivel de desequilibrio demográfico:
los clientes (hospitales) pueden tener distribuciones de edad o sexo
muy distintas entre sí, análogas al problema estudiado en este trabajo
pero a escala inter-institucional. Diseñar estrategias de balanceo
demográfico compatibles con el aprendizaje federado ---donde el balanceo
debe operar sobre los datos locales de cada cliente sin acceso a la
distribución global--- es una extensión natural de los métodos
desarrollados aquí y de alta relevancia para la implantación clínica
real de sistemas de predicción glucémica.

